\documentclass[../main.tex]{subfiles}

\begin{document}

\chapter{Daftar Referensi Tambahan dan Sumber Belajar}

Karya yang dikutip langsung dari bab-bab buku ini tercantum pada Daftar
Pustaka. Lampiran ini menambahkan rujukan pendalaman dan sumber belajar
eksternal yang dapat dipakai untuk memperluas materi.

\section{Referensi Tambahan}

\begin{description}
    \item[Stallings, W. (2017).] \textit{Cryptography and Network Security:
          Principles and Practice}, 7th ed. Pearson. Rujukan standar untuk
          fungsi hash, MAC, tanda tangan digital, PKI, dan protokol keamanan
          jaringan.
    \item[Menezes, A., van Oorschot, P., dan Vanstone, S. (1996).]
          \textit{Handbook of Applied Cryptography}. CRC Press. Rujukan
          ensiklopedis untuk landasan matematis dan algoritma kriptografi.
    \item[Katz, J. dan Lindell, Y. (2020).] \textit{Introduction to Modern
          Cryptography}, 3rd ed. CRC Press. Perlakuan formal untuk keamanan
          yang dapat dibuktikan (\textit{provable security}).
    \item[Ferguson, N. dan Schneier, B. (2003).] \textit{Practical
          Cryptography}. Wiley. Panduan penerapan kriptografi yang benar pada
          sistem nyata.
    \item[Paar, C. dan Pelzl, J. (2010).] \textit{Understanding Cryptography}.
          Springer. Pengantar yang menekankan contoh perhitungan.
\end{description}

\section{Sumber Belajar Eksternal}

\begin{sloppypar}
\begin{itemize}
    \item \textbf{Slide dan PDF kuliah IF4020 Kriptografi, Rinaldi Munir (ITB)}---
          \url{https://informatika.stei.itb.ac.id/~rinaldi.munir/Kriptografi/2023-2024/}.
          Bahan kuliah pengantar, cipher klasik, dan kriptografi modern yang
          menjadi rujukan utama silabus.
    \item \textbf{Handbook of Applied Cryptography, bab unduhan}---
          \url{http://cacr.uwaterloo.ca/hac/}. Bab 1 (ikhtisar) serta Bab 6--8
          (stream, block, dan kriptografi kunci-publik) tersedia bebas.
    \item \textbf{NIST Computer Security Resource Center}---\url{https://csrc.nist.gov}.
          Memuat standar \textit{Federal Information Processing Standards}
          (FIPS), termasuk FIPS 197 (AES), FIPS 180-4 (SHA), dan FIPS 186-5
          (tanda tangan digital).
    \item \textbf{IETF RFC 8446---TLS 1.3}---\url{https://www.rfc-editor.org/rfc/rfc8446}.
          Spesifikasi protokol TLS versi terbaru.
    \item \textbf{NIST Post-Quantum Cryptography}---\url{https://csrc.nist.gov/projects/post-quantum-cryptography}.
          Standar kriptografi yang tahan terhadap komputer kuantum.
    \item \textbf{Crypto I (CS255), Dan Boneh (Stanford)}---
          \url{https://crypto.stanford.edu/~dabo/courses/cs255/}. Materi kuliah
          kriptografi kunci-publik dan protokol keamanan.
    \item \textbf{Cryptohack}---\url{https://cryptohack.org}. Latihan interaktif
          kriptografi dari tingkat dasar sampai lanjut.
    \item \textbf{Massachusetts Institute of Technology OpenCourseWare}---
          \url{https://ocw.mit.edu}. Materi kuliah terbuka, antara lain kuliah
          kriptografi.
\end{itemize}
\end{sloppypar}

Sumber pada daftar di atas bersifat pilihan; mahasiswa dianjurkan
memprioritaskan Daftar Pustaka utama dan standar resmi (NIST, IETF) ketika
merujuk ke spesifikasi teknis.

\end{document}
